\renewcommand{\thesection}{\arabic{section}}\renewcommand{\theHsection}{chapter.\arabic{section}}\setcounter{section}{14}\section{函数列与函数项级数}
\noindent\begin{minipage}{\textwidth}\ProblemHeading{例题}{15.1（改写）}
求函数项级数 $\sum_{n=0}^{\infty}x^n$ 在实轴上的收敛点集与发散点集。
\end{minipage}\par\vfill
\noindent\begin{minipage}{\textwidth}\ProblemHeading{例题}{15.2（改写）}
设 $f_n(x)=x^n$（$0\le x\le1$）。求其逐点极限函数，并考察极限函数在 $x=1$ 处的连续性与可微性。
\end{minipage}\par\vfill
\noindent\begin{minipage}{\textwidth}\ProblemHeading{例题}{15.3（改写）}
把 $[0,1]$ 内的全部有理数排列为 $r_1,r_2,\ldots$，令 $S_n(x)$ 在 $\{r_1,\ldots,r_n\}$ 上等于 $1$，在其余点等于 $0$。证明每个 $S_n$ 在 $[0,1]$ 上 Riemann 可积，而其逐点极限不可积。
\end{minipage}\par\vfill
\newpage
\noindent\begin{minipage}{\textwidth}\ProblemHeading{例题}{15.4（改写）}
设 $S_n(x)=2n^2xe^{-n^2x^2}$（$0\le x\le1$）。求逐点极限 $S$，比较 $\lim_{n\to\infty}\int_0^1S_n(x)\,dx$ 与 $\int_0^1S(x)\,dx$。
\end{minipage}\par\vfill
\noindent\begin{minipage}{\textwidth}\ProblemHeading{例题}{15.5（改写）}
设 $S_n(x)=\sin(nx)/\sqrt n$（$x\in\mathbb R$）。求其极限函数，并说明函数列取极限与求导不一定可以交换。
\end{minipage}\par\vfill
\noindent\begin{minipage}{\textwidth}\ProblemHeading{习题}{15.6}
求下列函数项级数的收敛点集（各级数均从 $n=1$ 起求和）：
\SubQ{1}{$\displaystyle\sum\frac{n-1}{n+1}\left(\frac{x}{3x+1}\right)^n$。}
\SubQ{2}{$\displaystyle\sum ne^{-nx}$。}
\SubQ{3}{$\displaystyle\sum\left(\frac{x(x+n)}n\right)^n$。}
\end{minipage}\par\vfill
\newpage
\noindent\begin{minipage}{\textwidth}\ProblemHeading{习题}{15.7}
求下列函数项级数的收敛点集（各级数均从 $n=1$ 起求和）：
\SubQ{1}{$\displaystyle\sum\frac{x^n}{1+x^{2n}}$。}
\SubQ{2}{$\displaystyle\sum\frac{(n+x)^n}{n^{n+x}}$。}
\SubQ{3}{$\displaystyle\sum\frac{x^ny^n}{x^n+y^n}$，$x,y>0$。}
\end{minipage}\par\vfill
\noindent\begin{minipage}{\textwidth}\ProblemHeading{习题}{15.8}
求下列函数项级数的收敛点集（各级数均从 $n=1$ 起求和）：
\SubQ{1}{$\displaystyle\sum\frac{\ln(1+x^n)}{n^x}$，$x>0$。}
\SubQ{2}{$\displaystyle\sum(n^{1/n}-1)^x$，$x>0$。}
\end{minipage}\par\vfill
\noindent\begin{minipage}{\textwidth}\ProblemHeading{习题}{15.9}
设 $\sum_{n=1}^{\infty}u_n(x)$ 在 $[a,b]$ 上收敛，且每个 $u_n$ 非负、连续。证明其和函数在 $[a,b]$ 上能取到最小值。
\end{minipage}\par\vfill
\newpage
\noindent\begin{minipage}{\textwidth}\ProblemHeading{习题}{15.10}
设 $\sum_{n=1}^{\infty}u_n(x)$ 在 $[a,b]$ 上收敛，且每个 $u_n$ 非负、连续。其和函数是否一定能取到最大值？
\end{minipage}\par\vfill
\noindent\begin{minipage}{\textwidth}\ProblemHeading{习题}{15.11}
设区间 $I$ 为开区间或无界区间，$\sum_{n=1}^{\infty}u_n(x)$ 在 $I$ 上收敛，每个 $u_n$ 非负、连续。其和函数在 $I$ 上是否一定能取到最小值？
\end{minipage}\par\vfill
\noindent\begin{minipage}{\textwidth}\ProblemHeading{例题}{15.12（改写）}
考察下列函数列在 $(0,1)$ 上的逐点收敛性与一致收敛性，并比较它们的收敛情况：
\SubQ{1}{$f_n(x)=x^n$。}
\SubQ{2}{$g_n(x)=1/(n+x)$。}
\end{minipage}\par\vfill
\newpage
\noindent\begin{minipage}{\textwidth}\ProblemHeading{例题}{15.13}
考察 $f_n(x)=nx/(1+n^2x^2)$ 在 $(0,+\infty)$ 以及 $[\delta,+\infty)$（$\delta>0$）上的一致收敛性。
\end{minipage}\par\vfill
\noindent\begin{minipage}{\textwidth}\ProblemHeading{例题}{15.14}
考察 $f_n(x)=2n^2xe^{-n^2x^2}$ 在 $[0,1]$ 上的一致收敛性。
\end{minipage}\par\vfill
\noindent\begin{minipage}{\textwidth}\ProblemHeading{例题}{15.15（改写）}
令 $u_n(x)=1/(n+x)$（$n\ge1$，$x\in[0,+\infty)$）。证明 $u_n$ 一致趋于 $0$，但级数 $\sum_{n=1}^{\infty}u_n(x)$ 在每一点都发散。
\end{minipage}\par\vfill
\newpage
\noindent\begin{minipage}{\textwidth}\ProblemHeading{例题}{15.16}
考察级数 $\sum_{n=1}^{\infty}ne^{-nx}$ 在 $(0,+\infty)$ 上的一致收敛性。
\end{minipage}\par\vfill
\noindent\begin{minipage}{\textwidth}\ProblemHeading{例题}{15.17（改写）}
证明 $\sum_{n=1}^{\infty}\cos(nx)/n^2$ 在 $\mathbb R$ 上一致收敛。
\end{minipage}\par\vfill
\noindent\begin{minipage}{\textwidth}\ProblemHeading{例题}{15.18（改写）}
证明 $\sum_{n=1}^{\infty}x/(1+n^4x^2)$ 在 $\mathbb R$ 上一致收敛。
\end{minipage}\par\vfill
\newpage
\noindent\begin{minipage}{\textwidth}\ProblemHeading{例题}{15.19（改写）}
设 $\delta>0$。证明 $\sum_{n=1}^{\infty}ne^{-nx}$ 在 $[\delta,+\infty)$ 上一致收敛。
\end{minipage}\par\vfill
\noindent\begin{minipage}{\textwidth}\ProblemHeading{例题}{15.20（改写）}
证明函数列 $f_n(x)=x^n$ 在 $(0,1)$ 上一致有界。
\end{minipage}\par\vfill
\noindent\begin{minipage}{\textwidth}\ProblemHeading{例题}{15.21（改写）}
证明函数列 $f_n(x)=nx^n$ 在 $(0,1)$ 上逐点有界，但不一致有界。
\end{minipage}\par\vfill
\newpage
\noindent\begin{minipage}{\textwidth}\ProblemHeading{习题}{15.22}
函数项级数的 Abel 判别法假设 $\sum a_n(x)$ 在 $E$ 上一致收敛，$b_n(x)$ 对每个 $x$ 随 $n$ 单调且在 $E$ 上一致有界；Dirichlet 判别法则假设 $\sum_{k=1}^na_k(x)$ 一致有界、$b_n(x)$ 随 $n$ 单调且一致趋于 $0$。说明为什么不能直接照搬数项级数的证明，通过减去 $\lim_{n\to\infty}b_n(x)$ 来由后者推出前者。
\end{minipage}\par\vfill
\noindent\begin{minipage}{\textwidth}\ProblemHeading{例题}{15.23}
设 $0<\delta<\pi$。证明 $\sum_{n=1}^{\infty}\cos(nx)/n$ 与 $\sum_{n=1}^{\infty}\sin(nx)/n$ 在 $[\delta,2\pi-\delta]$ 上一致收敛。
\end{minipage}\par\vfill
\noindent\begin{minipage}{\textwidth}\ProblemHeading{例题}{15.24（改写）}
证明对每个 $x\in(0,\pi)$，级数 $\sum_{n=1}^{\infty}\cos(nx)/n$ 与 $\sum_{n=1}^{\infty}\sin(nx)/n$ 均不绝对收敛。
\end{minipage}\par\vfill
\newpage
\noindent\begin{minipage}{\textwidth}\ProblemHeading{例题}{15.25（改写）}
考察级数 $\sum_{n=0}^{\infty}(-1)^nx^n$ 在 $[0,1)$ 上的绝对收敛性与一致收敛性。
\end{minipage}\par\vfill
\noindent\begin{minipage}{\textwidth}\ProblemHeading{习题}{15.26（校订）}
考察下列函数列在指定区间上的一致收敛性：
\SubQ{1}{$f_n(x)=1/(1+nx)$，$x\in(0,+\infty)$。}
\SubQ{2}{$f_n(x)=1/(1+nx)$，$x\in(\lambda,+\infty)$，$\lambda>0$。}
\SubQ{3}{$f_n(x)=x^n/(1+x^n)$，$x\in[0,1-\lambda]$，$0<\lambda\le1$。}
\end{minipage}\par\vfill
\noindent\begin{minipage}{\textwidth}\ProblemHeading{习题}{15.27（校订）}
考察下列函数列在指定区间上的一致收敛性：
\SubQ{1}{$f_n(x)=x^n/(1+x^n)$，$x\in[1-\lambda,1+\lambda]$，$0<\lambda<2$。}
\SubQ{2}{$f_n(x)=x^n/(1+x^n)$，$x\in[1+\lambda,+\infty)$，$\lambda>0$。}
\SubQ{3}{$f_n(x)=e^{-(x-n)^2}$，$x\in(-l,l)$，$l>0$。}
\end{minipage}\par\vfill
\newpage
\noindent\begin{minipage}{\textwidth}\ProblemHeading{习题}{15.28（校订）}
考察下列函数列在指定区间上的一致收敛性：
\SubQ{1}{$f_n(x)=e^{-(x-n)^2}$，$x\in\mathbb R$。}
\end{minipage}\par\vfill
\noindent\begin{minipage}{\textwidth}\ProblemHeading{习题}{15.29（校订）}
考察下列级数在指定区间上的一致收敛性：
\SubQ{1}{$\displaystyle\sum_{n=1}^{\infty}\frac1{(x+n)(x+n+1)}$，$x>0$。}
\SubQ{2}{$\displaystyle\sum_{n=1}^{\infty}\frac{nx}{1+n^5x^2}$，$x\in\mathbb R$。}
\SubQ{3}{$\displaystyle\sum_{n=1}^{\infty}\frac{n^2}{\sqrt{n!}}(x^n+x^{-n})$，$x\in[-3,-1/3]\cup[1/3,3]$。}
\end{minipage}\par\vfill
\noindent\begin{minipage}{\textwidth}\ProblemHeading{习题}{15.30（校订）}
考察下列级数在指定区间上的一致收敛性：
\SubQ{1}{$\displaystyle\sum_{n=1}^{\infty}\frac{\sin((n+1/2)x)}{(n^4+x^4)^{1/3}}$，$x\in\mathbb R$。}
\SubQ{2}{$\displaystyle\sum_{n=2}^{\infty}\ln\left(1+\frac{x}{n\ln^2n}\right)$，$x\in(-l,l)$，$0<l\le2(\ln2)^2$。}
\SubQ{3}{$\displaystyle\sum_{n=2}^{\infty}\frac{(-1)^n}{n+\sin x}$，$x\in[0,2\pi]$。}
\end{minipage}\par\vfill
\newpage
\noindent\begin{minipage}{\textwidth}\ProblemHeading{习题}{15.31（校订）}
考察下列级数在指定区间上的一致收敛性：
\SubQ{1}{$\displaystyle\sum_{n=1}^{\infty}2^n\sin\frac1{3^nx}$，$x\in(0,+\infty)$。}
\end{minipage}\par\vfill
\noindent\begin{minipage}{\textwidth}\ProblemHeading{习题}{15.32}
若 $\sum_{n=1}^{\infty}a_n$ 收敛，证明 $\sum_{n=1}^{\infty}a_ne^{-nx}$ 在 $[0,+\infty)$ 上一致收敛。
\end{minipage}\par\vfill
\noindent\begin{minipage}{\textwidth}\ProblemHeading{习题}{15.33}
设每个 $u_n$ 在 $[a,b]$ 上单调，$\sum u_n(a)$ 与 $\sum u_n(b)$ 均绝对收敛。证明 $\sum u_n(x)$ 在 $[a,b]$ 上绝对且一致收敛。
\end{minipage}\par\vfill
\newpage
\noindent\begin{minipage}{\textwidth}\ProblemHeading{习题}{15.34}
证明 $\sum_{n=0}^{\infty}(-1)^nx^n(1-x)$ 在 $[0,1]$ 上绝对且一致收敛，而 $\sum_{n=0}^{\infty}x^n(1-x)$ 不一致收敛。
\end{minipage}\par\vfill
\noindent\begin{minipage}{\textwidth}\ProblemHeading{习题}{15.35}
在 $[0,1]$ 上令 $u_n(1/n)=1/n$，其余点 $u_n(x)=0$。证明 $\sum u_n(x)$ 一致收敛，但不存在收敛的正项级数 $\sum M_n$ 使 $|u_n(x)|\le M_n$ 对所有 $x,n$ 成立。
\end{minipage}\par\vfill
\noindent\begin{minipage}{\textwidth}\ProblemHeading{习题}{15.36}
设 $u_n$ 在 $[a,b]$ 上连续，$\sum u_n(x)$ 在 $[a,b)$ 上逐点收敛，而在 $b$ 处发散。证明该级数在 $[a,b)$ 上不一致收敛。
\end{minipage}\par\vfill
\newpage
\noindent\begin{minipage}{\textwidth}\ProblemHeading{习题}{15.37}
设 $a_n\ge0$。证明 $\sum_{n=1}^{\infty}a_n\cos(nx)$ 在 $\mathbb R$ 上一致收敛的充要条件为 $\sum a_n<+\infty$。
\end{minipage}\par\vfill
\noindent\begin{minipage}{\textwidth}\ProblemHeading{习题}{15.38}
证明 $\sum_{n=2}^{\infty}\cos(nx)/(n\ln n)$ 在 $(0,\pi]$ 上不一致收敛。
\end{minipage}\par\vfill
\noindent\begin{minipage}{\textwidth}\ProblemHeading{习题}{15.39}
设 $f_n,f$ 定义在有限闭区间 $[a,b]$ 上。若对每个 $x\in[a,b]$，都存在包含 $x$ 的开区间 $I_x$，使 $f_n$ 在 $I_x\cap[a,b]$ 上一致收敛于 $f$，证明 $f_n$ 在 $[a,b]$ 上一致收敛于 $f$。
\end{minipage}\par\vfill
\newpage
\noindent\begin{minipage}{\textwidth}\ProblemHeading{习题}{15.40（校订）}
设 $f_n(x)=xn^{-x}(\ln n)^\alpha$，$n\ge2$。证明它在 $[0,+\infty)$ 上一致收敛的充要条件是 $\alpha<1$。
\end{minipage}\par\vfill
\noindent\begin{minipage}{\textwidth}\ProblemHeading{习题}{15.41（校订）}
设 $f_1$ 在 $[a,b]$ 上 Riemann 可积，并递归定义 $f_{n+1}(x)=\int_a^xf_n(t)\,dt$。证明 $f_n$ 在 $[a,b]$ 上一致趋于 $0$。
\end{minipage}\par\vfill
\noindent\begin{minipage}{\textwidth}\ProblemHeading{习题}{15.42}
设 $\alpha>2$。证明 $\sum_{n=1}^{\infty}x^\alpha e^{-nx^2}$ 在 $[0,+\infty)$ 上一致收敛。
\end{minipage}\par\vfill
\newpage
\noindent\begin{minipage}{\textwidth}\ProblemHeading{习题}{15.43}
设 $\delta>0$。考察 $\sum_{n=1}^{\infty}\dfrac{nx}{(1+x)(1+2x)\cdots(1+nx)}$ 在 $[0,\delta]$ 和 $[\delta,+\infty)$ 上的一致收敛性。
\end{minipage}\par\vfill
\noindent\begin{minipage}{\textwidth}\ProblemHeading{习题}{15.44}
设 $f_n,g_n$ 在区间 $I$ 上一致收敛，且每个 $f_n,g_n$ 都在 $I$ 上有界。证明 $f_ng_n$ 在 $I$ 上一致收敛。
\end{minipage}\par\vfill
\noindent\begin{minipage}{\textwidth}\ProblemHeading{习题}{15.45}
若 $f_n,g_n$ 在区间 $I$ 上一致收敛，但不假设各函数有界，$f_ng_n$ 是否仍一定一致收敛？试举例说明。
\end{minipage}\par\vfill
\newpage
\noindent\begin{minipage}{\textwidth}\ProblemHeading{习题}{15.46}
设 $\sum u_n(x)$ 在 $[a,b]$ 上收敛，每个 $u_n$ 可导，且存在常数 $M$ 使所有 $x\in[a,b]$、$n\ge1$ 都满足 $|\sum_{k=1}^nu_k^{\prime}(x)|\le M$。证明该函数项级数在 $[a,b]$ 上一致收敛。
\end{minipage}\par\vfill
\noindent\begin{minipage}{\textwidth}\ProblemHeading{习题}{15.47}
设 $f$ 在 $0$ 的某邻域内二阶连续可导，$f(0)=0$，$0<f^{\prime}(0)<1$。令 $f_n$ 为 $f$ 的 $n$ 次复合。证明 $\sum_{n=1}^{\infty}f_n(x)$ 在 $0$ 的某邻域内一致收敛。
\end{minipage}\par\vfill
\noindent\begin{minipage}{\textwidth}\ProblemHeading{习题}{15.48}
证明 $\sum_{n=1}^{\infty}\dfrac{x^n\cos(nx)}{1+x+\cdots+x^{2n-1}}$ 在 $[0,1]$ 上一致收敛。
\end{minipage}\par\vfill
\newpage
\noindent\begin{minipage}{\textwidth}\ProblemHeading{习题}{15.49}
设 $a_n>0$ 且单调递减。证明 $\sum_{n=1}^{\infty}a_n\sin(nx)$ 在 $\mathbb R$ 上一致收敛的充要条件为 $na_n\to0$。
\end{minipage}\par\vfill
\noindent\begin{minipage}{\textwidth}\ProblemHeading{例题}{15.50（改写）}
证明 $f(x)=\sum_{n=1}^{\infty}\cos(nx)/n^2$ 在 $\mathbb R$ 上连续。
\end{minipage}\par\vfill
\noindent\begin{minipage}{\textwidth}\ProblemHeading{例题}{15.51}
设 $f(x)=\sum_{n=0}^{\infty}(x^n/3^n)\cos(n\pi x^2)$。计算 $\lim_{x\to1}f(x)$。
\end{minipage}\par\vfill
\newpage
\noindent\begin{minipage}{\textwidth}\ProblemHeading{例题}{15.52}
证明 $f(x)=\sum_{n=1}^{\infty}ne^{-nx}$ 在 $(0,+\infty)$ 上连续。
\end{minipage}\par\vfill
\noindent\begin{minipage}{\textwidth}\ProblemHeading{例题}{15.53}
证明 $f(x)=\sum_{n=2}^{\infty}(x/\ln n)^n$ 在 $\mathbb R$ 上连续。
\end{minipage}\par\vfill
\noindent\begin{minipage}{\textwidth}\ProblemHeading{例题}{15.54（改写）}
考察下列具体例子的一致收敛性，并说明为什么 Dini 定理中的有限闭区间条件不能任意去掉：
\SubQ{1}{函数列 $f_n(x)=x/n$，$x\in(0,+\infty)$。}
\SubQ{2}{函数项级数 $\sum_{n=0}^{\infty}x^n$，$x\in[0,1)$。}
\SubQ{3}{函数项级数 $\sum_{n=2}^{\infty}(x/\ln n)^n$，$x\in[0,+\infty)$。}
\end{minipage}\par\vfill
\newpage
\noindent\begin{minipage}{\textwidth}\ProblemHeading{例题}{15.55}
设 $f(x)=\sum_{n=1}^{\infty}\cos(nx)/n^2$。计算 $\int_0^\pi f(x)\,dx$。
\end{minipage}\par\vfill
\noindent\begin{minipage}{\textwidth}\ProblemHeading{例题}{15.56}
设 $f(x)=\sum_{n=1}^{\infty}\sin(nx)/n^4$。证明 $f\in C^2(\mathbb R)$，并求 $f^{\prime\prime}(x)$。
\end{minipage}\par\vfill
\noindent\begin{minipage}{\textwidth}\ProblemHeading{例题}{15.57（改写）}
已知当 $0<x<2\pi$ 时，$\sum_{n=1}^{\infty}\sin(nx)/n=(\pi-x)/2$。求此级数在整个实轴上的和函数及其不连续点，并说明连续函数级数的和函数不一定连续。
\end{minipage}\par\vfill
\newpage
\noindent\begin{minipage}{\textwidth}\ProblemHeading{习题}{15.58}
确定下列级数所定义函数的存在域，并讨论其连续性：
\SubQ{1}{$\displaystyle f(x)=\sum_{n=1}^{\infty}(1+1/n)^n$。}
\SubQ{2}{$\displaystyle f(x)=\sum_{n=1}^{\infty}\frac{x+(-1)^nn}{x^2+n^2}$。}
\end{minipage}\par\vfill
\noindent\begin{minipage}{\textwidth}\ProblemHeading{习题}{15.59}
设 $f(x)=\sum_{n=1}^{\infty}\cos(nx)/n^4$。证明 $f\in C^2(\mathbb R)$，并求 $f^{\prime\prime}$。
\end{minipage}\par\vfill
\noindent\begin{minipage}{\textwidth}\ProblemHeading{习题}{15.60}
证明 $\zeta(x)=\sum_{n=1}^{\infty}1/n^x$ 在 $(1,+\infty)$ 上连续且无穷次可导。
\end{minipage}\par\vfill
\newpage
\noindent\begin{minipage}{\textwidth}\ProblemHeading{习题}{15.61}
设 $f(x)=\sum_{n=1}^{\infty}ne^{-nx}$（$x>0$）。计算 $\int_{\ln2}^{\ln3}f(x)\,dx$。
\end{minipage}\par\vfill
\noindent\begin{minipage}{\textwidth}\ProblemHeading{习题}{15.62}
已知 $\sum_{n=1}^{\infty}(-1)^{n-1}/n=\ln2$，证明 $\lim_{x\to1}\sum_{n=1}^{\infty}(-1)^{n-1}/n^x=\ln2$。
\end{minipage}\par\vfill
\noindent\begin{minipage}{\textwidth}\ProblemHeading{习题}{15.63}
设 $E\subset\mathbb R$，$x_0$ 是 $E$ 的聚点（也允许 $x_0=\pm\infty$）。若 $\sum u_n(x)$ 在 $E$ 上一致收敛，且沿 $E$ 趋于 $x_0$ 时 $u_n(x)\to a_n$，证明 $\sum a_n$ 收敛且 $\lim_{x\to x_0,\,x\in E}\sum u_n(x)=\sum a_n$。
\end{minipage}\par\vfill
\newpage
\noindent\begin{minipage}{\textwidth}\ProblemHeading{习题}{15.64}
设 $f(x)=\sum_{n=1}^{\infty}\left(\dfrac{x}{1+2x}\right)^n\cos\dfrac{n\pi}x$。求 $\lim_{x\to1}f(x)$ 和 $\lim_{x\to+\infty}f(x)$。
\end{minipage}\par\vfill
\noindent\begin{minipage}{\textwidth}\ProblemHeading{习题}{15.65}
证明\[\lim_{x\to1}\sum_{n=1}^{\infty}\frac{x^n(1-x)}{n(1-x^{2n})}=\frac12\sum_{n=1}^{\infty}\frac1{n^2}.\]
\end{minipage}\par\vfill
\noindent\begin{minipage}{\textwidth}\ProblemHeading{习题}{15.66}
设 $\sum u_n(x)$ 在 $[a,b]$ 上逐点收敛，所有部分和在该区间上一致有界，每个 $u_n$ Riemann 可积。若对任意 $c\in(a,b)$ 及任意使区间非空的 $\delta>0$，该级数在 $[a,c-\delta]$ 与 $[c+\delta,b]$ 上一致收敛，证明和函数 Riemann 可积且可逐项积分。
\end{minipage}\par\vfill
\newpage
\noindent\begin{minipage}{\textwidth}\ProblemHeading{习题}{15.67}
设每个广义积分 $\int_a^{+\infty}u_n(t)\,dt$ 收敛，$\sum u_n(x)$ 在每个有限区间 $[a,b]$ 上一致收敛；记 $f_n(x)=\int_a^xu_n(t)\,dt$，并假设 $\sum f_n(x)$ 在 $[a,+\infty)$ 上一致收敛。证明\[\int_a^{+\infty}\sum_{n=1}^{\infty}u_n(x)\,dx=\sum_{n=1}^{\infty}\int_a^{+\infty}u_n(x)\,dx,\]且两边均收敛。
\end{minipage}\par\vfill
\noindent\begin{minipage}{\textwidth}\ProblemHeading{习题}{15.68}
设 $\sum u_n(x)$ 在 $x_0$ 的某邻域内收敛，每个 $u_n$ 在 $x_0$ 可导，且级数 $\sum\dfrac{u_n(x)-u_n(x_0)}{x-x_0}$ 在某去心邻域内一致收敛。证明和函数在 $x_0$ 可导，其导数等于 $\sum u_n^{\prime}(x_0)$。
\end{minipage}\par\vfill
\noindent\begin{minipage}{\textwidth}\ProblemHeading{习题}{15.69}
设 $a_1,a_2,\ldots$ 是 $(0,1)$ 内互异的点。证明 $f(x)=\sum_{n=1}^{\infty}|x-a_n|/2^n$ 在 $(0,1)$ 上连续，在每个 $a_n$ 处不可导，在 $(0,1)$ 内其余各点可导。
\end{minipage}\par\vfill
\newpage
\noindent\begin{minipage}{\textwidth}\ProblemHeading{习题}{15.70（校订）}
设 $f(x)=\sum_{n=0}^{\infty}1/(x+2^n)$（$x\ge0$）。证明：
\SubQ{1}{$f$ 在 $[0,+\infty)$ 上连续。}
\SubQ{2}{$\lim_{x\to+\infty}f(x)=0$。}
\SubQ{3}{对每个 $x>0$，有 $\displaystyle 0<f(x)-\frac{\ln(1+x)}{x\ln2}<\frac1{1+x}$。}
\end{minipage}\par\vfill
\noindent\begin{minipage}{\textwidth}\ProblemHeading{习题}{15.71}
证明 $\displaystyle\sum_{n=0}^{\infty}\frac{(n!)^2\,2^{n+1}}{(2n+1)!}=\pi$。
\end{minipage}\par\vfill
\noindent\begin{minipage}{\textwidth}\ProblemHeading{习题}{15.72}
构造一个可导函数，使它在每个有理点取有理数值，而它的导数在每个有理点取无理数值。
\end{minipage}\par\vfill
\newpage
\noindent\begin{minipage}{\textwidth}\ProblemHeading{例题}{15.73（改写）}
用 Cauchy--Hadamard 公式求下列幂级数的收敛半径：
\SubQ{1}{$\sum_{n=1}^{\infty}nx^n$。}
\SubQ{2}{$\sum_{n=0}^{\infty}x^n/n!$。}
\SubQ{3}{$\sum_{n=1}^{\infty}n^nx^n$。}
\end{minipage}\par\vfill
\noindent\begin{minipage}{\textwidth}\ProblemHeading{例题}{15.74}
求幂级数 $\sum_{k=0}^{\infty}2^kx^{2k}$ 的收敛半径。
\end{minipage}\par\vfill
\noindent\begin{minipage}{\textwidth}\ProblemHeading{例题}{15.75（改写）}
求下列幂级数的收敛半径，并分别判断在收敛区间两个端点处是绝对收敛、条件收敛还是发散：
\SubQ{1}{$\displaystyle\sum_{n=1}^{\infty}x^n/n$。}
\SubQ{2}{$\displaystyle\sum_{n=1}^{\infty}x^n/n^2$。}
\SubQ{3}{$\displaystyle\sum_{n=1}^{\infty}nx^n$。}
\end{minipage}\par\vfill
\newpage
\noindent\begin{minipage}{\textwidth}\ProblemHeading{例题}{15.76}
求幂级数 $\sum_{n=1}^{\infty}nx^n$ 的和函数。
\end{minipage}\par\vfill
\noindent\begin{minipage}{\textwidth}\ProblemHeading{例题}{15.77（改写）}
计算下列数项级数的和：
\SubQ{1}{$\displaystyle\sum_{n=1}^{\infty}n/2^n$。}
\SubQ{2}{$\displaystyle\sum_{n=1}^{\infty}n/3^n$。}
\end{minipage}\par\vfill
\noindent\begin{minipage}{\textwidth}\ProblemHeading{例题}{15.78}
把 $\ln(1+x)$ 与 $\arctan x$ 展开为以 $0$ 为中心的幂级数，并给出成立的范围。
\end{minipage}\par\vfill
\newpage
\noindent\begin{minipage}{\textwidth}\ProblemHeading{例题}{15.79}
计算下列级数的和：
\SubQ{1}{$\sum_{n=1}^{\infty}(-1)^{n-1}/n$。}
\SubQ{2}{$\sum_{n=0}^{\infty}(-1)^n/(2n+1)$。}
\end{minipage}\par\vfill
\noindent\begin{minipage}{\textwidth}\ProblemHeading{例题}{15.80}
计算 $\sum_{n=0}^{\infty}(-1)^n/(3n+1)$。
\end{minipage}\par\vfill
\noindent\begin{minipage}{\textwidth}\ProblemHeading{例题}{15.81（改写）}
设 $S(x)=\sum_{n=0}^{\infty}(-1)^nx^n$（$|x|<1$）。计算 $\lim_{x\to1^-}S(x)$，并判断 $\sum_{n=0}^{\infty}(-1)^n$ 的收敛性。由此说明：幂级数的和函数在收敛区间端点有有限单侧极限，并不能推出幂级数在该端点收敛。
\end{minipage}\par\vfill
\newpage
\noindent\begin{minipage}{\textwidth}\ProblemHeading{例题}{15.82（改写）}
设 $f(x)=\sum_{n=0}^{\infty}a_nx^n$ 的收敛半径为 $1$，$S_n=\sum_{k=0}^na_k$。证明在 $|x|<1$ 时：
\SubQ{1}{$\displaystyle\frac{f(x)}{1-x}=\sum_{n=0}^{\infty}S_nx^n$。}
\SubQ{2}{$\displaystyle\frac{f(x)}{(1-x)^2}=\sum_{n=0}^{\infty}(S_0+\cdots+S_n)x^n$。}
\end{minipage}\par\vfill
\noindent\begin{minipage}{\textwidth}\ProblemHeading{例题}{15.83}
证明当 $|x|<1$ 时，$\sum_{n=0}^{\infty}(n+1)(n+2)x^n=2/(1-x)^3$。
\end{minipage}\par\vfill
\noindent\begin{minipage}{\textwidth}\ProblemHeading{习题}{15.84}
求下列幂级数的收敛半径，并考察端点处的收敛性：
\SubQ{1}{$\sum_{n=1}^{\infty}(1+1/n)^{n^2}x^n$。}
\SubQ{2}{$\sum_{n=1}^{\infty}(1+1/2+\cdots+1/n)x^n$。}
\SubQ{3}{$\sum_{n=1}^{\infty}(a^n/n+b^n/n^2)x^n$，$a,b>0$。}
\end{minipage}\par\vfill
\newpage
\noindent\begin{minipage}{\textwidth}\ProblemHeading{习题}{15.85}
求下列幂级数的收敛半径，并考察端点处的收敛性：
\SubQ{1}{$\sum_{n=1}^{\infty}(n/e)^nx^n$。}
\end{minipage}\par\vfill
\noindent\begin{minipage}{\textwidth}\ProblemHeading{习题}{15.86}
求下列级数的收敛点集：
\SubQ{1}{$\displaystyle\sum_{n=0}^{\infty}\frac1{2n+1}\left(\frac{1-x}{1+x}\right)^n$。}
\SubQ{2}{$\displaystyle\sum_{n=1}^{\infty}(1+1/n)^{-n^2}e^{-nx}$。}
\SubQ{3}{$\displaystyle\sum_{n=1}^{\infty}x^{-n}\sin(\pi/2^n)$。}
\end{minipage}\par\vfill
\noindent\begin{minipage}{\textwidth}\ProblemHeading{习题}{15.87（校订）}
设幂级数 $\sum a_nx^n$ 与 $\sum b_nx^n$ 的收敛半径分别为 $R_1,R_2$。
\SubQ{1}{证明 $\sum(a_n+b_n)x^n$ 的收敛半径至少为 $\min(R_1,R_2)$。}
\SubQ{2}{若 $R_1R_2$ 有定义，证明 $\sum a_nb_nx^n$ 的收敛半径至少为 $R_1R_2$。}
\SubQ{3}{分别举例说明以上两个下界都可能严格小于相应的收敛半径。}
\end{minipage}\par\vfill
\newpage
\noindent\begin{minipage}{\textwidth}\ProblemHeading{习题}{15.88}
求下列幂级数在 $|x|<1$ 时的和函数：
\SubQ{1}{$\sum_{n=0}^{\infty}x^{2n+1}/(2n+1)$。}
\SubQ{2}{$\sum_{n=0}^{\infty}(-1)^{n-1}x^{2n+1}/(2n+1)$。}
\SubQ{3}{$\sum_{n=1}^{\infty}x^n/[n(n+1)]$。}
\end{minipage}\par\vfill
\noindent\begin{minipage}{\textwidth}\ProblemHeading{习题}{15.89}
证明当 $|x|<1$ 时：
\SubQ{1}{$\displaystyle\sum_{n=1}^{\infty}n^3x^n=\frac{x+4x^2+x^3}{(1-x)^4}$。}
\SubQ{2}{$\displaystyle\sum_{n=0}^{\infty}\frac{(n+1)(n+2)(n+3)}{3!}x^n=\frac1{(1-x)^4}$。}
\end{minipage}\par\vfill
\noindent\begin{minipage}{\textwidth}\ProblemHeading{习题}{15.90}
设 $\sum a_nx^n$ 的收敛半径为 $R$，$a_n$ 从某项起非零。令 $l=\liminf_{n\to\infty}|a_n/a_{n+1}|$，$L=\limsup_{n\to\infty}|a_n/a_{n+1}|$。证明 $l\le R\le L$。
\end{minipage}\par\vfill
\newpage
\noindent\begin{minipage}{\textwidth}\ProblemHeading{习题}{15.91}
设 $\sum_{n=0}^{\infty}a_nx^n$ 的收敛半径为 $1$，$a_n\ge0$，且 $\lim_{x\to1^-}\sum a_nx^n=A$ 为有限数。证明 $\sum a_n=A$。
\end{minipage}\par\vfill
\noindent\begin{minipage}{\textwidth}\ProblemHeading{习题}{15.92}
设 $\sum_{n=0}^{\infty}a_n=A$、$\sum_{n=0}^{\infty}b_n=B$，令 $c_n=\sum_{k=0}^na_kb_{n-k}$。若 $\sum c_n$ 收敛，利用幂级数乘法与 Abel 端点连续性定理证明 $\sum c_n=AB$。
\end{minipage}\par\vfill
\noindent\begin{minipage}{\textwidth}\ProblemHeading{习题}{15.93}
设正项级数 $\sum_{n=0}^{\infty}a_n$ 发散，且 $a_n/(a_0+\cdots+a_n)\to0$。证明 $\sum a_nx^n$ 的收敛半径为 $1$。
\end{minipage}\par\vfill
\newpage
\noindent\begin{minipage}{\textwidth}\ProblemHeading{习题}{15.94}
求级数 $1/2-1/5+1/8-1/11+\cdots$ 的和。
\end{minipage}\par\vfill
\noindent\begin{minipage}{\textwidth}\ProblemHeading{习题}{15.95}
证明 $\displaystyle1-\frac15+\frac19-\frac1{13}+\cdots=\frac{\pi+2\ln(\sqrt2+1)}{4\sqrt2}$。
\end{minipage}\par\vfill
\noindent\begin{minipage}{\textwidth}\ProblemHeading{习题}{15.96（校订）}
证明 $\displaystyle1+\sum_{n=1}^{\infty}\frac{(2n-1)!!}{(2n)!!}\frac1{(2n+1)^2}=\frac\pi2\ln2$。
\end{minipage}\par\vfill
\newpage
\noindent\begin{minipage}{\textwidth}\ProblemHeading{习题}{15.97（校订）}
设 $a_n\ge0$，幂级数 $\sum a_nx^n$ 的收敛半径满足 $0<R<+\infty$。证明 $\lim_{x\to R^-}\sum a_nx^n=\sum a_nR^n$（允许两边为 $+\infty$），并由此证明 $\sum_{n=1}^{\infty}1/n$ 发散。
\end{minipage}\par\vfill
\noindent\begin{minipage}{\textwidth}\ProblemHeading{例题}{15.98（改写）}
设 $f(x)=e^{-1/x^2}$（$x\ne0$），$f(0)=0$。已知 $f$ 无穷次可导且 $f^{(n)}(0)=0$（$n=0,1,\ldots$）。求 $f$ 在 $0$ 处的 Taylor 级数，并判断其和是否等于 $f(x)$。
\end{minipage}\par\vfill
\noindent\begin{minipage}{\textwidth}\ProblemHeading{例题}{15.99}
把 $e^x$ 展开为 Maclaurin 级数，并给出成立范围。
\end{minipage}\par\vfill
\newpage
\noindent\begin{minipage}{\textwidth}\ProblemHeading{例题}{15.100}
把 $\sin x$ 与 $\cos x$ 展开为 Maclaurin 级数，并给出成立范围。
\end{minipage}\par\vfill
\noindent\begin{minipage}{\textwidth}\ProblemHeading{例题}{15.101}
设 $\alpha\in\mathbb R$。把 $(1+x)^\alpha$ 展开为 Maclaurin 级数，并讨论成立范围及端点的收敛性。
\end{minipage}\par\vfill
\noindent\begin{minipage}{\textwidth}\ProblemHeading{例题}{15.102}
把 $1/[(1-x)(2-x)]$ 展开为 $x$ 的幂级数。
\end{minipage}\par\vfill
\newpage
\noindent\begin{minipage}{\textwidth}\ProblemHeading{例题}{15.103}
把 $\ln(1-x)/(1-x)$ 展开为 $x$ 的幂级数。
\end{minipage}\par\vfill
\noindent\begin{minipage}{\textwidth}\ProblemHeading{例题}{15.104}
把 $\ln((1+x)/(1-x))$ 展开为 $x$ 的幂级数。
\end{minipage}\par\vfill
\noindent\begin{minipage}{\textwidth}\ProblemHeading{习题}{15.105}
利用已知的初等函数展开式，将下列函数展开为 $x$ 的幂级数，并确定成立范围：
\SubQ{1}{$e^{x^2}$。}
\SubQ{2}{$\cos^2x$。}
\SubQ{3}{$x^{12}/(1-x)$。}
\end{minipage}\par\vfill
\newpage
\noindent\begin{minipage}{\textwidth}\ProblemHeading{习题}{15.106}
利用已知的初等函数展开式，将下列函数展开为 $x$ 的幂级数，并确定成立范围：
\SubQ{1}{$x/(1+x-2x^2)$。}
\SubQ{2}{$(1+x)e^{-x}$。}
\end{minipage}\par\vfill
\noindent\begin{minipage}{\textwidth}\ProblemHeading{习题}{15.107}
将下列函数展开为 $x$ 的幂级数，并确定成立范围：
\SubQ{1}{$(1+x)\ln(1+x)$。}
\SubQ{2}{$x\arctan x-\ln\sqrt{1+x^2}$。}
\SubQ{3}{$x/[(1-x)(1-x^2)]$。}
\end{minipage}\par\vfill
\noindent\begin{minipage}{\textwidth}\ProblemHeading{习题}{15.108}
将下列函数展开为 $x$ 的幂级数，并确定成立范围：
\SubQ{1}{$(1+x^2)\arctan x$。}
\end{minipage}\par\vfill
\newpage
\noindent\begin{minipage}{\textwidth}\ProblemHeading{习题}{15.109}
将下列函数展开为 $x$ 的幂级数，并确定成立范围：
\SubQ{1}{$\arcsin x$。}
\SubQ{2}{$\ln(x+\sqrt{1+x^2})$。}
\SubQ{3}{$\int_0^x(\sin t/t)\,dt$，被积函数在 $t=0$ 处取连续延拓值。}
\end{minipage}\par\vfill
\noindent\begin{minipage}{\textwidth}\ProblemHeading{习题}{15.110}
将 $\ln x$ 展开为 $(x-1)/(x+1)$ 的正整数次幂级数。
\end{minipage}\par\vfill
\noindent\begin{minipage}{\textwidth}\ProblemHeading{习题}{15.111}
将 $x/\sqrt{1+x}$ 展开为 $x/(1+x)$ 的正整数次幂级数。
\end{minipage}\par\vfill
\newpage
\noindent\begin{minipage}{\textwidth}\ProblemHeading{习题}{15.112}
设 $f(x)=\sum_{n=0}^{\infty}\sin(2^nx)/n!$。证明 $f$ 无穷次可导，但其在 $0$ 处的 Taylor 级数对一切 $x\ne0$ 都发散。
\end{minipage}\par\vfill
\noindent\begin{minipage}{\textwidth}\ProblemHeading{习题}{15.113}
设 $f$ 在 $[0,r]$ 上无穷次可导，且 $f$ 及其各阶导数均非负。证明 $f$ 在 $[0,r)$ 上等于其 Maclaurin 级数的和。
\end{minipage}\par\vfill
\noindent\begin{minipage}{\textwidth}\ProblemHeading{习题}{15.114}
证明 $\displaystyle1+\frac13-\frac15-\frac17+\frac19+\frac1{11}-\cdots=\frac{\sqrt2\,\pi}4$。
\end{minipage}\par\vfill
\newpage
\noindent\begin{minipage}{\textwidth}\ProblemHeading{习题}{15.115}
令 $S_n=\sum_{k=1}^{\infty}1/k^n$（$n\ge2$），$\gamma=\lim_{N\to\infty}(\sum_{k=1}^N1/k-\ln N)$。证明 $\sum_{n=2}^{\infty}(-1)^nS_n/n=\gamma$。
\end{minipage}\par\vfill
\noindent\begin{minipage}{\textwidth}\ProblemHeading{例题}{15.116（改写）}
对 $f\in C[0,1]$，定义 Bernstein 多项式 $B_n(f;x)=\sum_{k=0}^nf(k/n)\binom nk x^k(1-x)^{n-k}$。分别求下列函数的 $B_n(f;x)$：
\SubQ{1}{$f(x)=1$。}
\SubQ{2}{$f(x)=x$。}
\SubQ{3}{$f(x)=x^2$。}
\end{minipage}\par\vfill
\noindent\begin{minipage}{\textwidth}\ProblemHeading{例题}{15.117（改写）}
证明函数 $f(x)=1/x$ 不能在下列区间上用多项式一致逼近：
\SubQ{1}{$(0,1]$。}
\SubQ{2}{$[1,+\infty)$。}
\end{minipage}\par\vfill
\newpage
\noindent\begin{minipage}{\textwidth}\ProblemHeading{例题}{15.118（改写）}
定义 $B_n(f;x)=\sum_{k=0}^nf(k/n)\binom nk x^k(1-x)^{n-k}$。对 $f(x)=x^2$，计算 $\sup_{0\le x\le1}|B_n(f;x)-f(x)|$，并确定该误差趋于 $0$ 的阶。
\end{minipage}\par\vfill
\noindent\begin{minipage}{\textwidth}\ProblemHeading{习题}{15.119}
对 $f\in C[0,1]$，定义 Bernstein 多项式 $B_n(f;x)=\sum_{k=0}^nf(k/n)\binom nk x^k(1-x)^{n-k}$。证明：
\SubQ{1}{$B_n(f;0)=f(0)$，$B_n(f;1)=f(1)$。}
\SubQ{2}{若 $f\ge0$（或 $f\le0$），则 $B_n(f)\ge0$（或 $B_n(f)\le0$）。}
\SubQ{3}{若 $f$ 单调增加（或减少），则 $B_n(f)$ 也单调增加（或减少）。}
\end{minipage}\par\vfill
\noindent\begin{minipage}{\textwidth}\ProblemHeading{习题}{15.120}
对 $f\in C[0,1]$，定义 Bernstein 多项式 $B_n(f;x)=\sum_{k=0}^nf(k/n)\binom nk x^k(1-x)^{n-k}$。证明：
\SubQ{1}{若 $f$ 是凸函数，则 $B_n(f)$ 也是凸函数。}
\end{minipage}\par\vfill
\newpage
\noindent\begin{minipage}{\textwidth}\ProblemHeading{习题}{15.121}
设 $f\in C[a,b]$。证明：
\SubQ{1}{若 $\int_a^bf(x)x^n\,dx=0$ 对所有整数 $n\ge0$ 成立，则 $f\equiv0$。}
\SubQ{2}{若存在正整数 $N$，使 $\int_a^bf(x)x^n\,dx=0$ 对所有整数 $n\ge N$ 成立，则 $f\equiv0$。}
\end{minipage}\par\vfill
\noindent\begin{minipage}{\textwidth}\ProblemHeading{习题}{15.122}
设 $f\in C[0,1]$。若存在正整数 $k$，使 $\int_0^1f(x)x^{kn}\,dx=0$ 对每个整数 $n\ge1$ 成立，证明 $f\equiv0$。
\end{minipage}\par\vfill
\noindent\begin{minipage}{\textwidth}\ProblemHeading{习题}{15.123}
设 $f\in C[-1,1]$。证明：
\SubQ{1}{若所有奇次矩 $\int_{-1}^1f(x)x^{2n+1}\,dx=0$（$n\ge0$），则 $f$ 是偶函数。}
\SubQ{2}{若所有偶次矩 $\int_{-1}^1f(x)x^{2n}\,dx=0$（$n\ge0$），则 $f$ 是奇函数。}
\end{minipage}\par\vfill
\newpage
\noindent\begin{minipage}{\textwidth}\ProblemHeading{习题}{15.124}
证明：若 $f$ 能在整个实轴上用多项式一致逼近，则 $f$ 本身是多项式。
\end{minipage}\par\vfill
\noindent\begin{minipage}{\textwidth}\ProblemHeading{习题}{15.125}
设 $f\in C[1,+\infty)$ 且 $\lim_{x\to+\infty}f(x)=l$ 为有限数。证明对每个 $\varepsilon>0$，存在多项式 $P$，使所有 $x\ge1$ 都满足 $|f(x)-P(1/x)|<\varepsilon$。
\end{minipage}\par\vfill
\noindent\begin{minipage}{\textwidth}\ProblemHeading{习题}{15.126}
利用下列步骤给出 Weierstrass 逼近定理的另一证明：
\SubQ{1}{令 $c_n=\left(\int_{-1}^1(1-x^2)^n\,dx\right)^{-1}$，证明 $c_n<\sqrt n$（$n\ge1$）。}
\SubQ{2}{设 $f\in C[0,1]$，$f(0)=f(1)=0$，在区间外延拓为 $0$；令 $Q_n(t)=c_n(1-t^2)^n$，$P_n(x)=\int_{-1}^1f(x+t)Q_n(t)\,dt$。证明 $P_n$ 在 $[0,1]$ 上是多项式，并一致趋于 $f$。}
\SubQ{3}{去掉 $f(0)=f(1)=0$ 的条件，证明任意有限闭区间上的连续函数都能用多项式一致逼近。}
\end{minipage}\par\vfill
\newpage
\noindent\begin{minipage}{\textwidth}\ProblemHeading{习题}{15.127}
对 $f\in C[0,1]$，定义 Bernstein 多项式 $B_n(f;x)=\sum_{k=0}^nf(k/n)\binom nk x^k(1-x)^{n-k}$。若 $f\in C^1[0,1]$，证明 $B_n^{\prime}(f;x)$ 一致趋于 $f^{\prime}(x)$。
\end{minipage}\par\vfill
\noindent\begin{minipage}{\textwidth}\ProblemHeading{习题}{15.128}
设 $f\in C[a,b]$。证明存在多项式列 $P_n$ 在 $[a,b]$ 上一致趋于 $f$，且对每个 $x$，数列 $P_n(x)$ 单调递增。
\end{minipage}\par\vfill
\noindent\begin{minipage}{\textwidth}\ProblemHeading{例题}{15.129（改写）}
对实数 $\alpha$ 定义 $\binom\alpha0=1$，$\binom\alpha k=\alpha(\alpha-1)\cdots(\alpha-k+1)/k!$（$k\ge1$）。设 $\alpha,\beta\in\mathbb R$，$n\ge0$ 为整数。
\SubQ{1}{证明 $\displaystyle\sum_{k=0}^n\binom\alpha k\binom\beta{n-k}=\binom{\alpha+\beta}n$。}
\SubQ{2}{推出 $\displaystyle\sum_{k=0}^n\binom nk^2=\binom{2n}n$。}
\end{minipage}\par\vfill
\newpage
\noindent\begin{minipage}{\textwidth}\ProblemHeading{例题}{15.130}
设 $p,q,n$ 是非负整数。证明 $\sum_{k=0}^n\binom{p+k}p\binom{q+n-k}q=\binom{p+q+n+1}{p+q+1}$。
\end{minipage}\par\vfill
\noindent\begin{minipage}{\textwidth}\ProblemHeading{例题}{15.131（改写）}
设 $a_0=a_1=1$，$a_n=a_{n-1}+a_{n-2}$（$n\ge2$）。
\SubQ{1}{求数列 $a_n$ 的通项公式。}
\SubQ{2}{令 $\varphi=(1+\sqrt5)/2$，证明\[a_n=\begin{cases}\lfloor\varphi^{n+1}/\sqrt5\rfloor,&n\text{ 为奇数},\\\lfloor\varphi^{n+1}/\sqrt5\rfloor+1,&n\text{ 为偶数}.\end{cases}\]}
\end{minipage}\par\vfill
\noindent\begin{minipage}{\textwidth}\ProblemHeading{例题}{15.132}
求凸 $n$ 边形被 $n-3$ 条内部互不相交的对角线分成 $n-2$ 个三角形的不同分法数（$n\ge3$）。
\end{minipage}\par\vfill
\newpage
\noindent\begin{minipage}{\textwidth}\ProblemHeading{习题}{15.133}
设 $m\ge0,n\ge1$ 为整数。证明 $\sum_{k=0}^m\binom{k+n-1}{n-1}=\binom{m+n}n$。
\end{minipage}\par\vfill
\noindent\begin{minipage}{\textwidth}\ProblemHeading{习题}{15.134}
证明数列 $\binom{2n}n$（$n\ge0$）的母函数为 $(1-4x)^{-1/2}$。
\end{minipage}\par\vfill
\noindent\begin{minipage}{\textwidth}\ProblemHeading{习题}{15.135}
证明 $\sum_{k=0}^n\binom{2k}k\binom{2n-2k}{n-k}=2^{2n}$（$n\ge0$）。
\end{minipage}\par\vfill
\newpage
\noindent\begin{minipage}{\textwidth}\ProblemHeading{习题}{15.136}
证明 $\sum_{k=0}^nk\binom{2k}k\binom{2n-2k}{n-k}=n2^{2n-1}$（$n\ge1$）。
\end{minipage}\par\vfill
\noindent\begin{minipage}{\textwidth}\ProblemHeading{习题}{15.137}
设 $n,m\ge1$ 为整数，并约定当 $k>m$ 时 $\binom mk=0$。证明 $\sum_{k=0}^nk\binom nk\binom mk=m\binom{n+m-1}{n-1}$。
\end{minipage}\par\vfill
\noindent\begin{minipage}{\textwidth}\ProblemHeading{习题}{15.138}
求满足下列递推关系及初始条件的数列通项：
\SubQ{1}{$a_n=5a_{n-1}-6a_{n-2}$（$n\ge2$），$a_0=1,a_1=-2$。}
\SubQ{2}{$a_n=-a_{n-1}+16a_{n-2}-20a_{n-3}$（$n\ge3$），$a_0=0,a_1=1,a_2=-1$。}
\SubQ{3}{$a_n=2a_{n-1}-a_{n-2}+2^{n-2}$（$n\ge2$），$a_0=a_1=1$。}
\end{minipage}\par\vfill
\newpage
\noindent\begin{minipage}{\textwidth}\ProblemHeading{习题}{15.139}
证明 $\sum_{k=0}^n\dfrac1{1-2k}\binom{2k}k2^{2n-2k}=\binom{2n}n$（$n\ge0$）。
\end{minipage}\par\vfill
\noindent\begin{minipage}{\textwidth}\ProblemHeading{习题}{15.140}
计算 $\sum_{k=0}^{\lfloor n/2\rfloor}\binom{n-k}k$（$n\ge0$）。
\end{minipage}\par\vfill
\noindent\begin{minipage}{\textwidth}\ProblemHeading{习题}{15.141（校订）}
设 $\alpha\in\mathbb R$。计算 $\sum_{k=0}^{\lfloor n/2\rfloor}\binom{n-k}k\alpha^k$（$n\ge0$）。
\end{minipage}\par\vfill
\newpage
\noindent\begin{minipage}{\textwidth}\ProblemHeading{习题}{15.142}
对整数 $n\ge0$，利用参数和 $S_n(\alpha)=\sum_{k=0}^{\lfloor n/2\rfloor}\binom{n-k}k\alpha^k$ 的表达式，证明：
\SubQ{1}{$\displaystyle\sum_{k=0}^{\lfloor n/2\rfloor}2^k\binom{n-k}k=\frac{2^{n+1}+(-1)^n}3$。}
\SubQ{2}{$\displaystyle\sum_{k=0}^{\lfloor n/2\rfloor}(-1)^k\binom{n-k}k=\frac2{\sqrt3}\sin\frac{(n+1)\pi}3$。}
\end{minipage}\par\vfill
\noindent\begin{minipage}{\textwidth}\ProblemHeading{习题}{15.143}
证明 $\displaystyle\sum_{k=0}^{n-1}\left(\sum_{j=0}^k\binom nj\right)\left(\sum_{j=k+1}^n\binom nj\right)=\frac n2\binom{2n}n$（$n\ge1$）。
\end{minipage}\par\vfill
\noindent\begin{minipage}{\textwidth}\ProblemHeading{例题}{15.144（改写）}
令 $d(x)=\min_{k\in\mathbb Z}|x-k|$，并定义 $u_0(x)=d(x+1/2)$、$u_n(x)=4^{-n}d(4^nx)$（$n\ge1$）。证明 $f(x)=\sum_{n=0}^{\infty}u_n(x)$ 在实轴上处处连续、处处不可微。
\end{minipage}\par\vfill
\newpage
\noindent\begin{minipage}{\textwidth}\ProblemHeading{例题}{15.145（改写）}
在 $[0,1]$ 上定义 $\varphi(t)=0$（$0\le t\le1/3$）、$\varphi(t)=3t-1$（$1/3\le t\le2/3$）、$\varphi(t)=1$（$2/3\le t\le1$），并将其延拓为实轴上周期为 $2$ 的偶函数。令 $x(t)=\sum_{n=1}^{\infty}2^{-n}\varphi(3^{2n-2}t)$、$y(t)=\sum_{n=1}^{\infty}2^{-n}\varphi(3^{2n-1}t)$（$0\le t\le1$）。证明 $t\mapsto(x(t),y(t))$ 是连续曲线，其像为整个正方形 $[0,1]^2$。
\end{minipage}\par\vfill
\noindent\begin{minipage}{\textwidth}\ProblemHeading{习题}{15.146}
在 $[0,1]$ 上定义 $\varphi(t)=0$（$0\le t\le1/3$）、$\varphi(t)=3t-1$（$1/3\le t\le2/3$）、$\varphi(t)=1$（$2/3\le t\le1$），并将其延拓为实轴上周期为 $2$ 的偶函数。令\[\begin{aligned}x(t)&=\sum_{n=1}^{\infty}2^{-n}\varphi(3^{3n-3}t),\\y(t)&=\sum_{n=1}^{\infty}2^{-n}\varphi(3^{3n-2}t),\\z(t)&=\sum_{n=1}^{\infty}2^{-n}\varphi(3^{3n-1}t),\qquad0\le t\le1.\end{aligned}\]证明这条空间曲线连续，并填满整个立方体 $[0,1]^3$。
\end{minipage}\par\vfill
